\sgaStableTarget{sga6.e14.sec.5}\section*{5. Relations between @@M0@@ and @@M1@@}

\sgaStableTarget{sga6.e14.para.5-1}\textbf{5.1.} The example given in 4.8 does not rule out the possibility of finding, for an arbitrary scheme @@M2@@, a ring homomorphism
@@M3@@
such that, for every locally free Module @@M4@@ on @@M5@@, @@M6@@ takes the Chern class @@M7@@ of Exp. V to the cohomological Chern class of SGA 5 VII. Of course, @@M8@@ here denotes a prime distinct from the residue characteristics of @@M9@@. Since the ring @@M10@@ is generated by the @@M11@@, one sees in any case that @@M12@@ is uniquely determined by the requirement of compatibility with Chern classes.

\sgaStableTarget{sga6.e14.para.5-2}\textbf{5.2.} Using the Chern homomorphism
@@M13@@
and passing to the associated graded groups for the filtration of the first member and the degree filtration of the second, respectively, one obtains a homomorphism of graded groups (but not of rings!)
@@M14@@
satisfying the relation
@@M15@@
where @@M16@@ denotes the @@M17@@-adic Chern class. By division starting from the @@M18@@, one obtains a ring homomorphism:
@@M19@@
taking Chern classes to Chern classes, such that the image of the first member is contained in that of @@M20@@; consequently, the question of the existence of @@M21@@ has an affirmative answer upon replacing @@M22@@ by its quotient by its torsion subgroup. The first doubtful case in defining the components @@M23@@ of the hypothetical homomorphism @@M24@@ is that of @@M25@@; a priori, one succeeds only in defining @@M26@@. The answer to the existence question for @@M27@@ would be negative if, for example, one could find three invertible sheaves @@M28@@ on @@M29@@ satisfying the relation
@@M30@@
